
\documentclass{article}
\usepackage{colortbl}
\usepackage{makecell}
\usepackage{multirow}
\usepackage{supertabular}

\begin{document}

\newcounter{utterance}

\centering \large Interaction Transcript for game `codenames', experiment `unambiguous', episode 6 with reflexion{-}Qwen3.5{-}4B.
\vspace{24pt}

{ \footnotesize  \setcounter{utterance}{1}
\setlength{\tabcolsep}{0pt}
\begin{supertabular}{c@{$\;$}|p{.15\linewidth}@{}p{.15\linewidth}p{.15\linewidth}p{.15\linewidth}p{.15\linewidth}p{.15\linewidth}}
    \# & $\;$A & \multicolumn{4}{c}{Game Master} & $\;\:$B\\
    \hline

    \theutterance \stepcounter{utterance}  
    & & & \multicolumn{2}{p{0.3\linewidth}}{
        \cellcolor[rgb]{0.9,0.9,0.9}{
            \makecell[{{p{\linewidth}}}]{
                \texttt{\tiny{[GM$|$GM]}}
                {'hidden': {'team': ['jukebox', 'ketchup', 'inmate'], 'innocent': ['nuisance', 'tequila', 'decorator', 'stewardess', 'motorcade', 'pond'], 'opponent': [], 'assassin': ['titanium']}, 'revealed': {'team': {'team': [], 'innocent': [], 'opponent': [], 'assassin': []}, 'opponent': {'team': [], 'innocent': [], 'opponent': [], 'assassin': []}}}
            }
        }
    }
    & & \\ \\

    \theutterance \stepcounter{utterance}  
    & & \multicolumn{4}{p{0.6\linewidth}}{
        \cellcolor[rgb]{0.9,0.9,0.9}{
            \makecell[{{p{\linewidth}}}]{
                \texttt{\tiny{[P1$\langle$GM]}}
                \texttt{You have attempted this same task before and did not succeed. These are the plans you wrote after those attempts.} \\
\texttt{Trial \#1: New plan: I failed because I chose clues that were too broad or ambiguous, causing my partner to guess opponent words or distractors instead of my remaining team words. Specifically, "SODA" led to "concierge" (opponent) and "ketchup" (team), wasting a turn. "SUGAR" led to "toothpick" (team) but also "concierge" (opponent). I should have selected clues that are uniquely associated with only one or two of my remaining words (jukebox, ketchup, inmate) and have no semantic overlap with the opponent's remaining words (newcomer). For example, instead of "SODA" which relates to many things, I should have used a specific term like "MUSIC" for jukebox or "CONDIMENT" for ketchup, ensuring the partner doesn't accidentally guess "newcomer" or other distractors. I will focus on high{-}specificity, low{-}ambiguity clues for the final three words.} \\
\\ 
\texttt{Trial \#2: New plan: I failed because I only provided one target (ketchup) when I had three words (jukebox, ketchup, inmate) to protect. By leaving jukebox and inmate unclued, my partner was forced to guess distractors or opponent words to find them, wasting turns. Next time, I will provide a single{-}word clue that relates to all three of my team words simultaneously (e.g., "ENTERTAINMENT" for jukebox, "FOOD" for ketchup, and "PRISON" for inmate) to maximize the number of words guessed in one turn, ensuring I clear my entire list before the opponent can win.} \\
\\ 
\texttt{Trial \#3: New plan: I failed because I only clued "ketchup" with "FOOD", leaving "jukebox" and "inmate" unclued. This forced my partner to guess distractors or opponent words to find the remaining team words, wasting turns. Next time, I will select a single{-}word clue that is semantically related to all three of my team words (jukebox, ketchup, inmate) simultaneously, such as "AMUSEMENT" or "SIGHTSEEING", to maximize the number of words guessed in one turn and prevent the opponent from winning.} \\
\\ 
\\ 
\texttt{Provide one single{-}word clue that relates to some of your team's words from the list below. You can choose to how many and to which words your clue relates to.} \\
\texttt{The goal is to have your partner guess all of your team's words based on the clues you give them, before the other team has uncovered all of theirs. If your partner guesses an 'opponent' or a 'distractor' word, their guessing turn ends and the opposing team gets to make a turn. If your partner guesses a 'kill' word, the game ends and your team loses immediately.} \\
\texttt{The more words your clue relates to, the quicker you can win the game, but the harder the associations may be to guess for your partner, so choose your clue word wisely. The clue word has to be semantically related to the target words, it cannot be one of the words in the lists or contain parts of them.} \\
\\ 
\texttt{Always give your single{-}word clue and your comma{-}separated list of related target words in the following format and make your answers as short as possible, never include any other text than is required in this form:} \\
\\ 
\texttt{CLUE: <WORD>} \\
\texttt{TARGETS: <TARGETS>} \\
\\ 
\texttt{Your team words are: jukebox, ketchup, inmate.} \\
\texttt{Your opponent's team words are: .} \\
\texttt{Distractor words are: nuisance, tequila, decorator, stewardess, motorcade, pond.} \\
\texttt{The kill words are: titanium.} \\
            }
        }
    }
    & \\ \\

    \theutterance \stepcounter{utterance}  
    & \multicolumn{4}{p{0.6\linewidth}}{
        \cellcolor[rgb]{0.9,0.9,0.9}{
            \makecell[{{p{\linewidth}}}]{
                \texttt{\tiny{[P1$\rangle$GM]}}
                \texttt{CLUE: FOOD} \\
\texttt{TARGETS: ketchup} \\
            }
        }
    }
    & & \\ \\

    \theutterance \stepcounter{utterance}  
    & & & \multicolumn{2}{p{0.3\linewidth}}{
        \cellcolor[rgb]{0.9,0.9,0.9}{
            \makecell[{{p{\linewidth}}}]{
                \texttt{\tiny{[GM$|$GM]}}
                \texttt{food} \\
            }
        }
    }
    & & \\ \\

    \theutterance \stepcounter{utterance}  
    & & & \multicolumn{2}{p{0.3\linewidth}}{
        \cellcolor[rgb]{0.9,0.9,0.9}{
            \makecell[{{p{\linewidth}}}]{
                \texttt{\tiny{[GM$|$GM]}}
                ['ketchup']
            }
        }
    }
    & & \\ \\

    \theutterance \stepcounter{utterance}  
    & & & \multicolumn{2}{p{0.3\linewidth}}{
        \cellcolor[rgb]{0.9,0.9,0.9}{
            \makecell[{{p{\linewidth}}}]{
                \texttt{\tiny{[GM$|$GM]}}
                {'word': 'ketchup', 'assignment': 'team'}
            }
        }
    }
    & & \\ \\

    \theutterance \stepcounter{utterance}  
    & & & \multicolumn{2}{p{0.3\linewidth}}{
        \cellcolor[rgb]{0.9,0.9,0.9}{
            \makecell[{{p{\linewidth}}}]{
                \texttt{\tiny{[GM$|$GM]}}
                \texttt{opponent has won} \\
            }
        }
    }
    & & \\ \\

    \theutterance \stepcounter{utterance}  
    & \multicolumn{4}{p{0.6\linewidth}}{
        \cellcolor[rgb]{0.95,0.95,0.95}{
            \makecell[{{p{\linewidth}}}]{
                \texttt{\tiny{[P2$\rightarrow$P1]}}
                \texttt{Post{-}episode Reflexion (not sent to GM):} \\
\texttt{New plan: I failed because I only clued one word ("ketchup") while leaving two others ("jukebox", "inmate") unclued, forcing my partner to waste turns guessing distractors or opponent words to find the remaining targets. Next time, I will select a single{-}word clue that is semantically related to all three of my team words simultaneously (e.g., "AMUSEMENT" for jukebox, "FOOD" for ketchup, and "INMATE" for inmate) to maximize the number of words guessed in one turn and prevent the opponent from winning.} \\
            }
        }
    }
    & & \\ \\

\end{supertabular}
}

\end{document}
